% ============================================
% Extended Data Figures
% ============================================
\newpage
\subsection*{Extended Data Figures}

% Remove automatic "Figure N:" prefix since we label manually in captions
\captionsetup[figure]{name={},labelsep=none,labelformat=empty}

% Extended Data Fig. 1 - Taxonomy color map
\begin{figure}[H]
\centering
\includegraphics[width=0.9\textwidth]{figures/extended_data/ED_Fig1_combined.pdf}
\caption{\textbf{Extended Data Fig.~1. Taxonomic classification of bacterial isolates.} Table showing the 54 bacterial isolates used in the experiment, of which 4 pairs share identical ASV sequences. Each row shows the ASV identifier and its full taxonomic classification (Kingdom, Phylum, Class, Order, Family, Genus). Colors match those used in pie charts and composition figures throughout the manuscript.}
\label{fig:ed1}
\end{figure}

% Extended Data Fig. 2 - Sensitivity of outcome classification to metric choice
\begin{figure}[H]
\centering
\includegraphics[width=0.9\textwidth]{figures/extended_data/ED_Fig2_combined.pdf}
\caption{\textbf{Extended Data Fig.~2. Sensitivity of outcome classification to metric choice.} Stacked bar plot showing the distribution of coalescence outcomes (Dominance, Mixture, Restructuring) in Base medium using the primary similarity-map classification and four alternative metrics: Euclidean distance, Bray--Curtis dissimilarity, Jensen--Shannon divergence, and Jaccard index. Percentages and counts (n/total) are shown for each category. \rev{Dominance is recovered by abundance-weighted parental-similarity metrics, including the primary similarity-map classification, Euclidean distance, and Bray--Curtis dissimilarity, but not by Jaccard index or Jensen--Shannon divergence. This divergence reflects the different biological questions asked by each metric: abundance-weighted metrics quantify whether the coalesced community quantitatively resembles one parent more than the other, whereas Jaccard index measures presence/absence retention and Jensen--Shannon divergence measures differences across the full abundance distribution. Thus, Jaccard and Jensen--Shannon are best interpreted as complementary metric choices that emphasize different features of community composition.}}
\label{fig:ed2}
\end{figure}

% Extended Data Fig. 3 - Event-matched additive null model
\begin{figure}[H]
\centering
\includegraphics[width=\textwidth]{figures/extended_data/ED_Fig3_combined.pdf}
\caption{\textbf{Extended Data Fig.~3. Event-matched simple additive null model for Base medium coalescence outcomes.} \textbf{a}, Raw-count similarity map comparing event-matched additive null outcomes (orange open points; $n_{C,\mathrm{null}}=n_A+n_B$, constructed from each event's two parental raw 16S read-count vectors) with observed coalesced communities projected using the same raw-count representation (blue filled points); arrows show example null-to-observed pairs. \textbf{b}, Paired shifts in direction-independent parental asymmetry, $y = |2\mathrm{PDI}-1|$, from each additive null model outcome to the corresponding observed coalesced community for all 83 Base medium events; the dashed line marks the $y=0.5$ asymmetry boundary. \textbf{c}, Class transitions from additive null model outcomes (columns) to observed coalesced communities (rows). The additive null model is usually classified as Mixture, and the observed Base medium outcomes are shifted further toward parental asymmetry than their corresponding null model outcomes (77/83 events; one-sided paired Wilcoxon test, $p = 5.9 \times 10^{-14}$). See Supplementary Note~1 for full analysis and additional null models.}
\label{fig:ed3}
\end{figure}

% Extended Data Fig. 4 - Species number ablation
\begin{figure}[H]
\centering
\includegraphics[width=\textwidth]{figures/extended_data/ED_Fig4_combined.pdf}
\caption{\textbf{Extended Data Fig.~4. Effect of parental-community richness on coalescence outcomes.} \textbf{(a--e)} Similarity maps showing coalescence outcomes for simulated parental communities with initial richness of 4, 6, 12, 24, and 48 species at $\mu = 0.6$. The qualitative pattern of Dominance-dominated outcomes is consistent across parental-community richness values. \textbf{(f)} Stacked bar plots showing outcome fractions (Dominance, Mixture, Restructuring) across parental-community richness values at three interaction-strength parameter values ($\mu=0.3, 0.6, 0.8$). Dominance frequency remains relatively stable across parental-community richness values, while Mixture decreases and Restructuring increases with richer communities. Simulations used 200 independent replicates per condition with uniform interaction coefficients $\alpha_{ij} \sim U(0, 2\mu)$. See Supplementary Note~3 for full robustness analysis.}
\label{fig:ed4}
\end{figure}

% Extended Data Fig. 5 - Pairwise selection correlation vs interaction strength (simulation)
\begin{figure}[H]
\centering
\includegraphics[width=\textwidth]{figures/extended_data/ED_Fig5_combined.pdf}
\caption{\textbf{Extended Data Fig.~5. Pairwise selection correlation increases with the interaction-strength parameter in simulations.} \textbf{(a--c)} Pairwise selection correlation for species pairs from the same parental community (red, ``Same parental community'') versus cross-community pairs (blue, ``Cross-community'') at three representative interaction-strength parameter values. Gray horizontal lines indicate the random-selection baseline from an origin-shuffle null model. Individual dots show per-event correlations (100 stratified events displayed for legibility); squares with error bars show mean $\pm$ s.e.m. At low interaction-strength parameter value ($\mu = 0.3$), both within-community and cross-community pairs show positive correlations with a small difference ($\Delta = 0.05$). As $\mu$ increases ($\mu = 0.6$, $0.8$), within-community pairs maintain positive correlation while cross-community pairs become increasingly negative, consistent with stronger origin-correlated species fates. \textbf{(d)} Mean pairwise selection correlation across all simulated coalescence events as a function of the interaction-strength parameter $\mu$. Species pairs from the same parental community (red circles) show positive correlation that increases with $\mu$, while cross-community pairs (blue squares) show negative correlation that becomes more negative with $\mu$. Error bars represent event-level s.e.m.\ across $n = 1{,}200$ coalescence events per interaction-strength parameter value generated from 200 independent replicates, with six pairwise coalescence events per replicate.}
\label{fig:ed5}
\end{figure}

% Extended Data Fig. 6 - Pairwise selection correlation (experimental)
\begin{figure}[H]
\centering
\includegraphics[width=\textwidth]{figures/extended_data/ED_Fig6_combined.pdf}
\caption{\textbf{Extended Data Fig.~6. Pairwise selection correlation in experimental coalescence across nutrient conditions.} Mean pairwise selection correlation for species pairs from the same parental community (red) versus cross-community pairs (blue). Gray horizontal line indicates the random-selection baseline from an origin-shuffle null model. Individual dots show per-event correlations; squares with error bars show mean $\pm$ s.e.m. \textbf{(a)} In Nutr$-$ medium, where invasion resistance and net interaction intensity are low, within-community and cross-community pairs show no significant difference in selection correlation, consistent with species-level dynamics and no detectable community-level selection. \textbf{(b)} In Base medium, intermediate differences emerge. \textbf{(c)} In Nutr$+$ medium, where invasion resistance and environmental feedbacks are stronger, within-community species pairs exhibit significantly higher selection correlation than cross-community pairs ($p < 0.001$), consistent with stronger origin-correlated species fates associated with community-level selection.}
\label{fig:ed6}
\end{figure}

% Extended Data Fig. 7 - Assembly effect simulation
\begin{figure}[H]
\centering
\includegraphics[width=0.95\textwidth]{figures/extended_data/ED_Fig7_combined.pdf}
\caption{\textbf{Extended Data Fig.~7. Assembly history promotes Dominance in simulations.} Stacked bar plots comparing outcome distributions between coalescence of pre-assembled parental communities (top row) and direct assembly from the mixed species pool (bottom row) across five interaction-strength parameter values ($\mu = 0.2, 0.4, 0.6, 0.8, 1.0$). Percentages and counts (n/total) are shown for each category. Coalescence consistently produces higher Dominance and lower Restructuring compared to direct assembly, supporting the interpretation that assembly history contributes to origin-correlated persistence.}
\label{fig:ed7}
\end{figure}

% Extended Data Fig. 8 - pH difference vs coalescence outcome
\begin{figure}[H]
\centering
\includegraphics[width=0.85\textwidth]{figures/extended_data/ED_Fig8_combined.pdf}
\caption{\textbf{Extended Data Fig.~8. Parental-community pH difference is associated with Dominance direction in Nutr$+$.} In the stricter acid--alkaline subset used for this PDI-regression analysis, where one parental community is acidic (pH $<$ 6.5) and the other alkaline (pH $>$ 7.5), the pH difference is associated with which parental community dominates in Nutr$+$ but not in Base medium. \textbf{(a)} Base medium ($n = 41$; $R^2 = 0.02$, $p = 0.42$, n.s.) and \textbf{(b)} Nutr$+$ medium ($n = 32$; $R^2 = 0.29$, $p = 0.002$). X-axis: pH difference (alkaline parental community pH $-$ acidic parental community pH). Y-axis: PDI transformed so larger values indicate greater contribution from the acidic parental community. Red shaded region indicates acidic community dominance (PDI $>$ 0.6); blue shaded region indicates alkaline community dominance (PDI $<$ 0.4). The relationship is significant only in Nutr$+$, consistent with pH-mediated environmental feedback contributing to the top-down regime rather than fully explaining the Dominance pattern. A broader acid--alkaline-pair control is shown in Supplementary Fig.~37.}
\label{fig:ed8}
\end{figure}
